% ═══════════════════════════════════════════════════════════════
% MUSTERLÖSUNG: Reaktion von Metallen mit Sauerstoff
% Thema 2.4 | Chemie Klasse 8 | DS 1-2
% ═══════════════════════════════════════════════════════════════

\documentclass[11pt, a4paper]{article}

\usepackage[utf8]{inputenc}
\usepackage[T1]{fontenc}
\usepackage[ngerman]{babel}

% --- SCHRIFTEN ---
\usepackage{helvet}
\renewcommand{\familydefault}{\sfdefault}

\usepackage{microtype}
\usepackage[margin=1.5cm, top=2cm, bottom=1.5cm]{geometry}
\usepackage{enumitem}
\usepackage{tabularx}
\usepackage{booktabs}
\usepackage{array}
\usepackage{multicol}
\usepackage{tikz}
\usetikzlibrary{calc}

\usepackage{amsmath, amssymb}
\usepackage{siunitx}
\sisetup{locale = DE, per-mode = symbol}

\usepackage{fancyhdr}
\usepackage{lastpage}
\usepackage{needspace}

\usepackage{xcolor}
\usepackage{tcolorbox}
\tcbuselibrary{skins}

\usepackage{qrcode}

% ═══════════════════════════════════════════════════════════════
% FARBEN
% ═══════════════════════════════════════════════════════════════

\definecolor{chemie}{HTML}{2a7a4b}
\definecolor{hellgrau}{HTML}{F5F5F5}
\definecolor{mittelgrau}{HTML}{888888}
\definecolor{dunkelgrau}{HTML}{444444}
\definecolor{rahmen}{HTML}{CCCCCC}
\definecolor{loesung}{HTML}{C62828}  % Rot für Lösungen

\colorlet{fachfarbe}{chemie}

% ═══════════════════════════════════════════════════════════════
% VARIABLEN
% ═══════════════════════════════════════════════════════════════

\newcommand{\fach}{Chemie}
\newcommand{\thema}{Reaktion von Metallen mit Sauerstoff}
\newcommand{\klasse}{8}
\newcommand{\maxpunkte}{40}

% Differenzierung anzeigen (Kl. 6-9)
\newif\ifzeigesternchen\zeigesternchentrue

% ═══════════════════════════════════════════════════════════════
% KOPF- UND FUSSZEILE
% ═══════════════════════════════════════════════════════════════

\pagestyle{fancy}
\fancyhf{}
\renewcommand{\headrulewidth}{0.4pt}
\renewcommand{\footrulewidth}{0pt}
\lhead{\small \fach{} Klasse \klasse{} | \thema}
\rhead{\small\textcolor{loesung}{\textbf{MUSTERLÖSUNG}}}
\cfoot{\small\textcolor{mittelgrau}{Seite \thepage{} von \pageref{LastPage}}}

% ═══════════════════════════════════════════════════════════════
% BOXEN
% ═══════════════════════════════════════════════════════════════

\newtcolorbox{merkkasten}[1][]{
    colback=hellgrau,
    colframe=hellgrau,
    boxrule=0pt,
    arc=0pt,
    left=8pt, right=8pt, top=8pt, bottom=6pt,
    borderline north={2pt}{0pt}{fachfarbe},
    before skip=6pt, after skip=6pt,
    #1
}

\newtcolorbox{formelkasten}{
    colback=white,
    colframe=fachfarbe,
    boxrule=1pt,
    arc=0pt,
    left=8pt, right=8pt, top=6pt, bottom=6pt,
    before skip=4pt, after skip=4pt
}

% ═══════════════════════════════════════════════════════════════
% BEFEHLE
% ═══════════════════════════════════════════════════════════════

\newcommand{\luecke}[1]{\underline{\hspace{#1}}}
\newcommand{\lsg}[1]{\textcolor{loesung}{\textbf{#1}}}  % Lösung rot
\newcommand{\punktebox}[1]{\hfill\small\textcolor{mittelgrau}{[\_\_/#1]}}

\newcommand{\stern}[1]{%
    \ifzeigesternchen%
        \textsuperscript{\textcolor{fachfarbe}{(#1)}}%
    \fi%
}

\newcommand{\aufgabe}[2]{%
    \needspace{4\baselineskip}%
    \vspace{10pt}%
    \noindent\textbf{\textcolor{fachfarbe}{Aufgabe #1}} #2%
    \vspace{4pt}%
    \nopagebreak[4]%
}

\newcommand{\operator}[1]{\textbf{\textcolor{fachfarbe}{#1}}}

\newlist{teil}{enumerate}{2}
\setlist[teil,1]{label=\alph*), leftmargin=1.8em, itemsep=3pt, parsep=0pt, topsep=3pt}
\setlist[teil,2]{label=\roman*), leftmargin=1.8em, itemsep=2pt}

\newcommand{\antwortzeilen}[1]{%
    \par\vspace{4pt}%
    \foreach \n in {1,...,#1}{%
        \noindent\rule{\linewidth}{0.3pt}\vspace{6pt}%
    }%
}

\setlength{\parindent}{0pt}
\setlength{\parskip}{4pt}
\setlength{\columnsep}{1cm}

% ═══════════════════════════════════════════════════════════════
% DOKUMENT
% ═══════════════════════════════════════════════════════════════

\begin{document}

% ═══════════════════════════════════════════════════════════════
% HEADER
% ═══════════════════════════════════════════════════════════════

\begin{center}
    {\Large\bfseries\textcolor{fachfarbe}{\thema}}\\[0.2em]
    {\small Musterlösung | \fach{} Klasse \klasse{} | Thema 2.4}
\end{center}
\vspace{-0.3em}
\hrule
\vspace{0.6em}

% ═══════════════════════════════════════════════════════════════
% KASTEN 1: MERKSÄTZE
% ═══════════════════════════════════════════════════════════════

\begin{merkkasten}
\textbf{Merksatz 1 – Oxidation:}\\[0.3em]
\lsg{Die Oxidation ist die chemische Reaktion eines Stoffes mit Sauerstoff.}\\[0.3em]
\lsg{Bei der Oxidation entsteht ein Oxid und es wird Energie freigesetzt.}

\vspace{0.8em}

\textbf{Merksatz 2 – Oxid:}\\[0.3em]
Ein \lsg{Oxid} ist eine chemische \lsg{Verbindung} aus einem Element und \lsg{Sauerstoff}.
\end{merkkasten}

% ═══════════════════════════════════════════════════════════════
% AUFGABEN
% ═══════════════════════════════════════════════════════════════

\aufgabe{1}{\stern{*} Begriffe zuordnen} \punktebox{4}

\operator{Ordne} die Begriffe den Definitionen zu.

\begin{center}
\begin{tabular}{|p{3cm}|p{10cm}|}
\hline
\textbf{Begriff} & \textbf{Definition} \\
\hline
\lsg{Oxidation} & Die chemische Reaktion eines Stoffes mit Sauerstoff. \\
\hline
\lsg{Element} & Ein Stoff, der aus nur einer Atomsorte besteht. \\
\hline
\lsg{Verbindung} & Ein Stoff, der aus verschiedenen Atomsorten in festem Verhältnis besteht. \\
\hline
\lsg{Oxid} & Eine chemische Verbindung aus einem Element und Sauerstoff. \\
\hline
\end{tabular}
\end{center}

% ---

\aufgabe{2}{\stern{*} Wortgleichungen formulieren} \punktebox{6}

\operator{Ergänze} die Wortgleichungen für die Reaktion von Metallen mit Sauerstoff.

\begin{teil}
    \item Magnesium + Sauerstoff $\longrightarrow$ \lsg{Magnesiumoxid}
    \item Eisen + Sauerstoff $\longrightarrow$ \lsg{Eisenoxid}
    \item \lsg{Kupfer} + Sauerstoff $\longrightarrow$ Kupferoxid
    \item Aluminium + Sauerstoff $\longrightarrow$ \lsg{Aluminiumoxid}
    \item Zink + \lsg{Sauerstoff} $\longrightarrow$ Zinkoxid
    \item \lsg{Calcium} + \lsg{Sauerstoff} $\longrightarrow$ Calciumoxid
\end{teil}

% ---

\aufgabe{3}{\stern{**} Formelgleichungen aufstellen} \punktebox{6}

\operator{Schreibe} zu den Wortgleichungen die Formelgleichungen.

\begin{teil}
    \item Magnesium + Sauerstoff $\longrightarrow$ Magnesiumoxid

    \vspace{0.3em}
    \lsg{2 Mg} + \lsg{O$_2$} $\longrightarrow$ \lsg{2 MgO}

    \item Aluminium + Sauerstoff $\longrightarrow$ Aluminiumoxid

    \vspace{0.3em}
    \lsg{4 Al} + \lsg{3 O$_2$} $\longrightarrow$ \lsg{2 Al$_2$O$_3$}

    \item Eisen + Sauerstoff $\longrightarrow$ Eisenoxid (Fe$_2$O$_3$)

    \vspace{0.3em}
    \lsg{4 Fe} + \lsg{3 O$_2$} $\longrightarrow$ \lsg{2 Fe$_2$O$_3$}
\end{teil}

% ---

\newpage

\aufgabe{4}{\stern{**} Beobachtung und Protokoll} \punktebox{6}

\textbf{a)} Bei der Verbrennung von Magnesium kann man beobachten:

\begin{center}
\begin{tabular}{|p{4cm}|p{9cm}|}
\hline
\textbf{Vor der Reaktion} & \lsg{Silbrig glänzendes Magnesiumband} \\
\hline
\textbf{Während der Reaktion} & \lsg{Sehr helle, weiße Flamme; starke Hitze} \\
\hline
\textbf{Nach der Reaktion} & \lsg{Weißes, sprödes Pulver (Magnesiumoxid)} \\
\hline
\textbf{Energieumwandlung} & \lsg{Chemische Energie → Licht + Wärme (exotherm)} \\
\hline
\end{tabular}
\end{center}

\vspace{0.3em}

\textbf{b)} Flammenfarben der Metalle (aus Simulation):

\begin{center}
\begin{tabular}{|c|c|c|}
\hline
\textbf{Metall} & \textbf{Flammenfarbe} & \textbf{Oxid-Farbe} \\
\hline
Magnesium (Mg) & \lsg{weiß (gleißend)} & \lsg{weiß} \\
\hline
Eisen (Fe) & \lsg{orange (Glut)} & \lsg{braun/rostrot} \\
\hline
Aluminium (Al) & \lsg{silber-weiß} & \lsg{grau} \\
\hline
Kupfer (Cu) & \lsg{grün-blau} & \lsg{schwarz} \\
\hline
\end{tabular}
\end{center}

% ---

\aufgabe{5}{\stern{**} Massenerhaltung} \punktebox{8}

\begin{teil}
    \item \operator{Erkläre} das Gesetz von der Erhaltung der Masse.

    \lsg{Bei einer chemischen Reaktion bleibt die Gesamtmasse aller beteiligten Stoffe erhalten. Es gehen keine Atome verloren und es entstehen keine neuen.}

    \begin{formelkasten}
    \textbf{Gesetz von der Erhaltung der Masse:}\\[0.3em]
    \lsg{Die Gesamtmasse der Ausgangsstoffe ist gleich der Gesamtmasse der Reaktionsprodukte.}
    \end{formelkasten}

    \item \operator{Berechne:} 12\,g Magnesium reagieren mit 8\,g Sauerstoff.

    Masse des Magnesiumoxids: \lsg{12\,g + 8\,g = \textbf{20\,g}}
\end{teil}

% ---

\aufgabe{6}{\stern{**} Korrosion im Alltag} \punktebox{4}

Rost ist ein bekanntes Oxidationsprodukt im Alltag.

\begin{teil}
    \item \operator{Nenne} zwei weitere Beispiele für Korrosion im Alltag.

    \lsg{1. Grünspan auf Kupferdächern (Kupferoxid/Kupfercarbonat)}

    \lsg{2. Angelaufenes Silberbesteck (Silbersulfid, nicht Oxidation)}

    \item \operator{Erkläre}, warum Aluminium nicht rostet, obwohl es mit Sauerstoff reagiert.

    \lsg{Aluminium bildet sofort eine dünne, fest haftende Oxidschicht, die das Metall schützt (Passivierung).}
\end{teil}

% ---

\aufgabe{7}{\stern{***} Vergleich: Magnesium und Aluminium} \punktebox{6}

\operator{Vergleiche} die Oxidation von Magnesium und Aluminium mithilfe der Simulation.

\begin{center}
\begin{tabular}{|p{4.5cm}|p{4.5cm}|p{4.5cm}|}
\hline
\textbf{Kriterium} & \textbf{Magnesium} & \textbf{Aluminium} \\
\hline
Oxidschicht-Struktur & \lsg{porös (durchlässig)} & \lsg{dicht (undurchlässig)} \\[0.5em]
\hline
Reaktion stoppt? & \lsg{Nein, vollständige Oxidation} & \lsg{Ja, nur Oberfläche} \\[0.5em]
\hline
Fachbegriff für das Verhalten & \lsg{Kettenreaktion / Durchoxidation} & \lsg{Passivierung} \\[0.5em]
\hline
\end{tabular}
\end{center}

\operator{Erkläre} den Unterschied in 1--2 Sätzen:

\lsg{Bei Magnesium ist die Oxidschicht porös, sodass Sauerstoff weiter zum Metall vordringen kann – die Reaktion läuft vollständig ab (exotherme Kettenreaktion).}

\lsg{Bei Aluminium bildet sich eine dichte Oxidschicht (Al$_2$O$_3$), die weiteren Sauerstoff blockiert – die Reaktion stoppt an der Oberfläche (Passivierung).}

% ═══════════════════════════════════════════════════════════════
% WERTIGKEITSTABELLE
% ═══════════════════════════════════════════════════════════════

\vspace{0.5em}

\begin{tcolorbox}[colback=hellgrau, colframe=hellgrau, boxrule=0pt, arc=0pt, left=6pt, right=6pt, top=4pt, bottom=4pt]
\small
\textbf{Wertigkeiten:} Mg(II), Al(III), Fe(II/III), Cu(I/II), Zn(II), Ca(II), O(II)
\end{tcolorbox}

% ═══════════════════════════════════════════════════════════════
% PUNKTE
% ═══════════════════════════════════════════════════════════════

\vfill
\hrule
\vspace{0.4em}

\hfill\textbf{Punkte:} \luecke{1.2cm} / \maxpunkte

\end{document}
